% Section: Method
% Outline (for review):
% - Definitions: probe prompting; cross-prompt activation signatures and metrics
% - Node selection via cumulative influence threshold
% - Preprocessing: functional vs semantic, target-token windowing and directionality
% - Decision rules: Semantic(Dictionary/Concept), Relationship, Say-X; duplicate resolution; stability
% - Naming policy and blacklist
% - Metrics: Neuronpedia Replacement/Completeness; behavioral coherence metrics
% - Implementation notes: scripts, UI, rate limits, checkpoints, seeds

\section{Method}
\label{sec:method}

\paragraph{Probe prompting}
Given an attribution graph for a seed prompt and target logit, we elicit candidate concepts via an instructed LLM. For each accepted concept, we create \emph{probe prompts} that preserve syntactic structure while varying semantic content. We then measure feature activations on each probe using the Neuronpedia API, forming a set of \emph{cross-prompt activation signatures}.

\subsection{Core Definitions}

\paragraph{Cross-prompt activation signature}
A feature-specific behavioral profile measured across probe prompts. For each (feature, probe) pair, we record: (i) cosine similarity to seed activation pattern (how similar is activation to original prompt?), (ii) robust z-score relative to baseline distribution (how anomalous is activation?), (iii) peak token position and identity (which token activates most strongly?), (iv) activation density (fraction of tokens above 75th percentile), and (v) sparsity ratio $(\text{peak}-\text{mean})/\text{peak}$ (how concentrated is activation?). 

Aggregating across probes yields summary statistics: (i) peak-token consistency (fraction of probes peaking on same token type), (ii) number of distinct peak tokens (diversity of activation), (iii) functional vs semantic peak percentage (does feature peak on `is'/`the' or on content words?), and (iv) confidence scores when peaking on semantic or functional tokens (reliability of peak type). These signatures capture whether a feature responds consistently to a concept (high consistency, few peaks), covers multiple related tokens (moderate consistency, several semantic peaks), or activates diffusely across contexts (low sparsity, many peaks).

\paragraph{Functional vs semantic tokens}
We label each token in the prompt as either \emph{functional} (bridging/syntactic role: `is', `the', `of', `and', etc.) or \emph{semantic} (content-bearing: entity names, concepts, attributes). This distinction is crucial for Say-X detection: features that consistently peak on functional tokens but map to nearby semantic targets (via directionality rules) are classified as output promoters. Complete functional vocabulary and mapping rules provided in Appendix~\ref{sec:functional-vocab}.

\paragraph{Replacement vs Completeness (Neuronpedia graph metrics)}
\emph{Replacement} measures the fraction of end-to-end influence from input embeddings to target logit that routes through pinned features in the subgraph (high Replacement = features carry most causal pathways). \emph{Completeness} measures the fraction of incoming edge influence to each node that is explained by upstream features/embeddings (high Completeness = few unexplained inputs, subgraph is self-contained). Both are computed via attribution path-tracing in replacement models~\cite{ameisen2025circuittracing}. Our subgraphs intentionally trade Replacement ($\sim$0.54) for Completeness ($\sim$0.83) and legibility: unpinned low-influence nodes are treated as error, reducing raw coverage but improving interpretability.

\paragraph{Node selection}
We select candidate feature nodes with a cumulative influence threshold on the attribution graph. The threshold is chosen via a simple UI that visualizes node influence; selected nodes form the universe for cross-prompt measurement.

\paragraph{Cross-prompt signature metrics}
For each feature and probe we record: (i) cosine similarity to the seed activation pattern, (ii) robust z-score relative to baseline, (iii) peak token, (iv) activation density (fraction above the 75th percentile), and (v) sparsity ratio $(\text{peak}-\text{mean})/\text{peak}$. Aggregating across probes yields: (i) peak consistency on concept tokens, (ii) number of distinct peaks, (iii) functional vs semantic peak percentage, and (iv) confidence scores when peaking on semantic or functional tokens.

\paragraph{Preprocessing}
Tokens are labeled as \emph{semantic} (content-bearing) or \emph{functional} (bridging tokens such as ``is'' and ``the''). When a feature peaks on a functional token, we identify the nearest semantic \emph{target token} within a configurable window (default: 7 tokens), using directionality rules (``is'' $\rightarrow$ forward, ``of'' $\rightarrow$ backward). See Appendix~\ref{sec:functional-vocab} for complete vocabulary and mapping rules.

\subsection{Decision Rules and Classification Logic}

We use transparent, rule-based thresholds rather than learned classifiers to ensure interpretability of the grouping process itself---a critical property for an interpretability tool. Learned cluster algorithms (k-means, DBSCAN, neural groupers) lack explicit decision criteria, making it difficult to understand \emph{why} a feature was assigned to a category. Rule-based classification provides:

\begin{itemize}
  \item \textbf{Transparency:} Each assignment can be traced to specific threshold crossings
  \item \textbf{Debuggability:} Misclassifications can be diagnosed by inspecting which condition failed
  \item \textbf{Editability:} Users can adjust thresholds for domain-specific needs
  \item \textbf{Communicability:} Rules can be stated in natural language
\end{itemize}

Thresholds were chosen via iterative refinement on held-out examples, balancing recall (capturing relevant features in each category) against precision (avoiding false groupings). While hand-tuned, they represent codified expert judgment. Classification proceeds hierarchically:

\begin{itemize}
  \item \textbf{Semantic (Dictionary):} stable token detector for a specific wordform or name; peaks on the same token across contexts. Peak consistency $\ge 0.80$ on a single token type; few distinct peaks ($\le 1$).
  \item \textbf{Semantic (Concept):} category detector that fires on a set of semantically related content tokens (synonyms, inflections, near-neighbors).Lower layers ($\le 3$) or high semantic confidence; can cover related content tokens.
  \item \textbf{Relationship:} sentence-spanning activations that accentuate conceptually linked tokens. Diffuse activations (median sparsity $< 0.45$) across relation phrases; often early-layer.
  \item \textbf{Say-X:} output-promotion feature that pushes the model to emit a specific token or short phrase, typically via functional tokens. Functional dominance ($\ge 50\%$), high functional consistency ($\ge 0.90$), and layer $\ge 7$.
\end{itemize}

\paragraph{Conflict resolution}
Features may satisfy multiple category conditions (e.g., high peak consistency for both semantic and functional tokens). We resolve conflicts via weighted alignment score combining: (i) peak consistency (weight 0.4), (ii) category-specific confidence (weight 0.3), (iii) layer prior (weight 0.2), and (iv) sparsity consistency (weight 0.1). Feature is assigned to highest-scoring category. Duplicate prevention ensures each feature belongs to exactly one supernode.

\paragraph{Cross-prompt stability requirement}
Features must satisfy conditions on $\ge$60\% of probes to avoid spurious groupings based on single-probe noise. Features failing this criterion are marked ``ungrouped'' and excluded from subgraph (typically 5--10\% of candidate features).

\paragraph{Naming policy}
Semantic nodes are named by their strongest semantic peak token; Say-X nodes are named by the target token discovered from functional peaks; Relationship nodes use aggregated activation over an extended vocabulary (format ``(token) related''). We support a configurable blacklist for uninformative words.

\paragraph{Naming conventions and vocabulary}
Functional token vocabulary (English prompts): `is', `was', `are', `were', `be', `been', `being' (copulas); `the', `a', `an' (articles); `of', `in', `on', `at', `to', `for', `with', `by', `from', `as' (prepositions); `and', `or', `but' (conjunctions); `that', `which', `who', `whose', `where', `when' (relative pronouns). Full list in Appendix~\ref{sec:functional-vocab}.

Target-token mapping uses $\pm$5 token window with directionality: Forward: `is', `was', `are' $\rightarrow$ look forward for nearest semantic token; Backward: `of', `'s' (possessive) $\rightarrow$ look backward; Bidirectional: `,' $\rightarrow$ nearest semantic token (either direction).

\paragraph{Cross-lingual note}
Non-English prompts require adapted functional vocabularies. We encountered a notable failure when applying the English pipeline to the French prompt ``le contraire de `petit' est'' using English probe prompts. The model produced incoherent activations with no valid supernode grouping, due to language mismatch and absence of French functional tokens. Performance improved after adding French probe prompts and expanding the blacklist, but the current prototype remains unreliable in multilingual contexts without explicit cross-lingual adaptation.

\paragraph{Metrics}
We report Neuronpedia \emph{Replacement} (fraction of end-to-end influence through features) and \emph{Completeness} (fraction of incoming edge influence explained by features/tokens) on full graphs and subgraphs. For interpretability, we use: (i) peak token consistency, (ii) within-cluster activation similarity, and (iii) sparsity consistency. We compare against cosine-only and layer-adjacency baselines.

\subsection{Implementation and Availability}

Our pipeline is implemented in Python 3.9+. Core dependencies: PyTorch for model inference, OpenAI/Anthropic APIs for probe generation, numpy/pandas for data processing, requests for Neuronpedia API, streamlit for interactive UI (see \texttt{requirements.txt} for complete list). The codebase is publicly available at \url{https://github.com/peppinob-ol/attribution-graph-probing} under GPL v3 license with:

\begin{itemize}
  \item \textbf{Stand-alone scripts} (\texttt{scripts/}) for programmatic execution and batch processing
  \item \textbf{Interactive Streamlit UI} (\texttt{eda/}) for exploratory analysis with visual threshold selection
  \item \textbf{Automatic rate limiting} with exponential backoff retry logic (respects Neuronpedia API limits)
  \item \textbf{Per-feature checkpoints} enabling resume from partial runs without re-querying API
  \item \textbf{Deterministic pipeline} ensuring reproducibility (no random sampling or initialization)
  \item \textbf{Example data} (\texttt{examples\_data/}) with outputs for all five circuits in this paper
\end{itemize}

We provide an interactive demo at \url{https://huggingface.co/spaces/Peppinob/attribution-graph-probing} enabling researchers to:

\begin{enumerate}
  \item Upload custom attribution graphs (Neuronpedia JSON format)
  \item Adjust cumulative influence threshold $\tau$ via visual curve interface
  \item Generate probe prompts with LLM assistance (OpenAI/Anthropic API)
  \item Review and edit proposed concepts before probe synthesis
  \item Download concept-aligned subgraphs for import to Neuronpedia workspace
\end{enumerate}

The demo requires no local installation and processes typical graphs (50--200 features) in 2--5 minutes. Users need only a Neuronpedia API key (free, available at neuronpedia.org). Activation measurement requires approximately 10--15 minutes on a standard GPU (L4) for 5 probes $\times$ 40 features. See Appendix~\ref{sec:repro} for detailed reproducibility instructions and CLI usage.


